\documentclass[9pt, conference, a4paper]{IEEEtran}
\IEEEoverridecommandlockouts
% The preceding line is only needed to identify funding in the first footnote. If that is unneeded, please comment it out.

% --- FONT AYARI (Arial / Helvetica) ----------------------
\usepackage[T1]{fontenc}
\usepackage[scaled]{helvet}
\renewcommand{\familydefault}{\sfdefault}
% --------------------------------------------------------------------------

\usepackage{cite}
\usepackage{amsmath,amssymb,amsfonts}
\usepackage{algorithmic}
\usepackage{graphicx}
\usepackage{textcomp}
\usepackage{xcolor}
\usepackage{array}
\def\BibTeX{{\rm B\kern-.05em{\sc i\kern-.025em b}\kern-.08em
    T\kern-.1667em\lower.7ex\hbox{E}\kern-.125emX}}
\begin{document}

\title{Scene-Routed Experts: Using Cheap Scene Labels to Improve Multi-Label Aerial Image Classification on AID\\
{\footnotesize This project proposal is submitted for the CS415 course at Sabanci University}
}

\author{
\IEEEauthorblockN{Yağız Çetin}
\IEEEauthorblockA{\textit{Computer Science and Engineering} \\
\textit{Sabanci University}\\
Istanbul, Turkey \\
yagiz.cetin@sabanciuniv.edu}
\\[0.5cm]
\IEEEauthorblockN{Erkan Ulaş Tepe}
\IEEEauthorblockA{\textit{Computer Science and Engineering} \\
\textit{Sabanci University}\\
Istanbul, Turkey \\
ulas.tepe@sabanciuniv.edu}
\and
\IEEEauthorblockN{Öykü Demirel}
\IEEEauthorblockA{\textit{Computer Science and Engineering} \\
\textit{Sabanci University}\\
Istanbul, Turkey \\
oyku.demirel@sabanciuniv.edu}
\\[0.5cm]
\IEEEauthorblockN{Bilge Kağan Yılmaz}
\IEEEauthorblockA{\textit{Computer Science and Engineering} \\
\textit{Sabanci University}\\
Istanbul, Turkey \\
bilge.yilmaz@sabanciuniv.edu}
}

\maketitle

\begin{abstract}
Multi-label aerial image classification predicts which objects appear in an aerial image. In the AID dataset, only 3,000 images have object labels, but all 10,000 images have a cheap scene label such as port or airport. We propose a model in which the predicted scene decides how a set of small expert classifiers is combined, so that the scene changes both which objects are expected and how the image is interpreted. We will test whether this improves object recognition and how many object labels it can save.
\end{abstract}

\begin{IEEEkeywords}
multi-label classification, aerial images, scene context, mixture of experts
\end{IEEEkeywords}

\section{Problem Definition and Motivation}
Most work on aerial images gives each image a single scene label such as \textit{river} or \textit{airport} \cite{xia2017aid}, but one image usually contains many objects. Two \textit{river} images may look completely different: one shows only water and trees, the other also buildings, roads and cars. Multi-label classification describes this content by predicting which of a fixed set of objects appear in the image, without the cost of pixel-level or bounding-box annotation \cite{hua2020relation}.

The task is hard because objects vary in size and appearance, some objects tend to appear together (ships and water), labels are very imbalanced (\textit{tree} appears in 2,406 images, \textit{airplane} in 99), and object labels are scarce: in the AID dataset only 3,000 of 10,000 images have them. However, \emph{all} 10,000 images have a cheap scene label, and knowing that an image shows a port makes ships and docks much more likely. Our question is: \textit{can these scene labels help a model recognize objects, and what is the best way to use them?}

\section{Related Work}
Xia et al. \cite{xia2017aid} built AID with 10,000 images in 30 scene classes. Hua et al. labeled 3,000 of them with 17 objects and modeled how labels depend on each other with a relation network \cite{hua2020relation}; later work used transformers \cite{wu2022smat}. None of these methods uses the scene labels or the 7,000 extra images.

Using the scene to recognize objects is an old idea \cite{torralba2010forest}. SALGL \cite{zhu2023salgl} uses scene-specific label statistics for natural images, but it has to discover the scenes by itself, while AID gives us real scene labels. Since scene labels are used only during training, our setting is an example of learning with privileged information \cite{vapnik2009lupi}. Our model is a mixture of experts \cite{jacobs1991moe}, where a router decides how much each expert contributes; recent work shows that guiding the router with labels can help \cite{zeng2026med}, but this has not been tried for multi-label aerial images.

\section{Dataset and Proposed Method}
\textbf{Dataset.} We use the multi-label AID dataset \cite{hua2020relation} from Hugging Face: 3,000 images of $600 \times 600$ pixels from all 30 scenes. We drop \textit{mobile home}, which appears in only two images, and keep 16 labels. Since this release has no scene labels, we will find each image in the original AID by image matching, which also gives us the 7,000 scene-only images. Inside each scene we split the 3,000 images into 2,160 training, 240 validation and 600 test images, close to the 80/20 split of \cite{hua2020relation}. Scene labels of test images are never shown to the model.

\textbf{Method (SR-MoE).} The same pattern can mean different things in different scenes: a flat blue area is \textit{water} in a port but may be a \textit{court} in a stadium. So the scene should change not only \emph{which} objects are likely, but also \emph{how} the image is read. On top of a pretrained ResNet-50 \cite{he2016resnet}, a scene head predicts the scene, and this prediction decides how a set of small, scene-specific expert classifiers is mixed to produce the 16 object scores. The experts share most of their weights, so each one can still predict all labels from little data. The scene head is trained on all 10,000 images, so the 7,000 images without object labels also help. To avoid shortcuts such as ``airport, so there must be an airplane'', we sometimes turn the scene input off during training.

\textbf{Experiments.} We compare (B0) a plain ResNet-50, (B1) a model that only learns the scene as a second task, (B2) a model that uses the scene only as a prior on object frequencies, and (B3) our SR-MoE. We also train B3 with a router that learns without scene labels, and repeat B0 and B3 with 10\%, 25\% and 50\% of the object labels.

\textbf{Hypotheses.} \textbf{H1:} Scene labels improve object recognition (B1 $>$ B0). \textbf{H2:} Changing the classifier per scene beats using the scene only as a prior (B3 $>$ B2), mostly for scene-specific objects such as airplanes, ships and docks. \textbf{H3:} Real scene labels make a better router than no labels. \textbf{H4:} With scene labels, half of the object labels are enough to match the baseline trained with all of them.

\section{Performance Metrics and Success Criteria}
Our main metric is \textbf{mAP} (mean average precision over the 16 labels). We also report example-based F1, used by \cite{hua2020relation}, per-label AP and scene accuracy. To check shortcuts, we count predicted objects that are typical for the scene but absent from the image. Results are averaged over 3 runs; a gain counts only if it exceeds the run-to-run variation. Our targets are given in Table~\ref{tab:criteria}; the project is successful if the baseline target and at least two hypotheses are met.

\begin{table}[htbp]
\caption{Success criteria}
\begin{center}
\begin{tabular}{|>{\raggedright\arraybackslash}p{5.6cm}|>{\raggedright\arraybackslash}p{1.8cm}|}
\hline
\textbf{Criterion} & \textbf{Target} \\
\hline
B0 example-F1 vs. 86.23 in \cite{hua2020relation} & within $\pm 1.5$ \\
\hline
H1: mAP gain of B1 over B0 & $\geq +1.0$ \\
\hline
H2: mAP gain of B3 over B2 & $\geq +1.0$ \\
\hline
H2: AP gain of B3 over B0, scene-specific & $\geq +3.0$ \\
\hline
H3: mAP gain, labeled vs. unlabeled router & $\geq +1.0$ \\
\hline
H4: B3 (50\% labels) minus B0 (100\%) & $\geq 0$ \\
\hline
Shortcut errors of B3 vs. B2 & not higher \\
\hline
\end{tabular}
\label{tab:criteria}
\end{center}
\end{table}

\section{Project Responsibilities}
\textbf{Erkan Ulaş Tepe:} data preparation, scene label matching, splits, baseline B0. \textbf{Öykü Demirel:} models B1 and B2, router without scene labels (H3). \textbf{Bilge Kağan Yılmaz:} reduced-label experiments (H4) and shortcut analysis. \textbf{Yağız Çetin:} SR-MoE (B3), evaluation code and visualizations. \textbf{All:} literature review, report and presentation.

\begin{thebibliography}{00}
\bibitem{xia2017aid} G.-S. Xia, J. Hu, F. Hu, B. Shi, X. Bai, Y. Zhong, L. Zhang, and X. Lu, ``AID: A benchmark data set for performance evaluation of aerial scene classification,'' \textit{IEEE Trans. Geosci. Remote Sens.}, vol. 55, no. 7, pp. 3965--3981, 2017.
\bibitem{hua2020relation} Y. Hua, L. Mou, and X. X. Zhu, ``Relation network for multilabel aerial image classification,'' \textit{IEEE Trans. Geosci. Remote Sens.}, vol. 58, no. 7, pp. 4558--4572, 2020.
\bibitem{wu2022smat} H. Wu, C. Xu, and H. Liu, ``S-MAT: Semantic-driven masked attention transformer for multi-label aerial image classification,'' \textit{Sensors}, vol. 22, no. 14, art. 5433, 2022.
\bibitem{torralba2010forest} A. Torralba, K. P. Murphy, and W. T. Freeman, ``Using the forest to see the trees: Exploiting context for visual object detection and localization,'' \textit{Commun. ACM}, vol. 53, no. 3, pp. 107--114, 2010.
\bibitem{zhu2023salgl} X. Zhu, J. Liu, W. Liu, J. Ge, B. Liu, and J. Cao, ``Scene-aware label graph learning for multi-label image classification,'' in \textit{Proc. ICCV}, 2023, pp. 1473--1482.
\bibitem{vapnik2009lupi} V. Vapnik and A. Vashist, ``A new learning paradigm: Learning using privileged information,'' \textit{Neural Networks}, vol. 22, no. 5--6, pp. 544--557, 2009.
\bibitem{jacobs1991moe} R. A. Jacobs, M. I. Jordan, S. J. Nowlan, and G. E. Hinton, ``Adaptive mixtures of local experts,'' \textit{Neural Computation}, vol. 3, no. 1, pp. 79--87, 1991.
\bibitem{zeng2026med} Z. Zeng, D. Sridhar, and N. Vasconcelos, ``MED-DSLC: Multi-expert-domain classification via domain supervision and logit calibration,'' arXiv:2607.10985, 2026.
\bibitem{he2016resnet} K. He, X. Zhang, S. Ren, and J. Sun, ``Deep residual learning for image recognition,'' in \textit{Proc. CVPR}, 2016, pp. 770--778.
\end{thebibliography}

\end{document}
